\section{Method}
\label{sec:method}

\subsection{Pipeline Overview}
\label{sec:method:pipeline}

\aris{} transforms a natural-language object description into a set of
high-quality rendered images through six stages (Figure~\ref{fig:pipeline}).
The pipeline is fully automated: no human intervention is required between
the initial concept description and the final quality-approved renders.

\paragraph{Stage 1: Text expansion.}
Given a seed concept (e.g., \textit{``wooden chair''}),
a language model expands it into a detailed text-to-image (T2I) prompt
that specifies material, color, geometry, and surface texture.
This step increases prompt specificity, improving the quality and diversity
of downstream generated images.

\paragraph{Stage 2: Object image generation.}
A T2I model (Qwen-Image-2512~\cite{PLACEHOLDER_qwen_image_2512_verify})
generates a $1024 \times 1024$ reference image of the object
against a clean background, using the expanded prompt.

\paragraph{Stage 2.5: Foreground segmentation.}
SAM3~\cite{PLACEHOLDER_sam3_verify} segments the object foreground from the background,
producing an RGBA image with transparent background.
This removes any incidental background content while preserving the object with high fidelity.
Critically, SAM3 operates on the object region only, so the resulting 3D reconstruction
(Stage 3) does not contain background geometry such as a floor plane,
eliminating the need for post-hoc mesh cleaning.

\paragraph{Stage 3: Single-image 3D reconstruction.}
Hunyuan3D-2.1~\cite{PLACEHOLDER_hunyuan3d_verify} reconstructs a textured 3D mesh (GLB format)
from the RGBA image using a feedforward single-image reconstruction model.
The resulting mesh includes physically-based rendering (PBR) texture maps,
enabling realistic rendering under varied lighting conditions.

\paragraph{Stage 4: Scene-aware Blender rendering.}
The reconstructed mesh is inserted into a real Blender scene and rendered using
Cycles path tracing.
Section~\ref{sec:method:blender} describes this stage in detail.

\paragraph{Stage 5.5--5.6: VLM review and rendering evolution.}
Rendered images are reviewed by a VLM (Qwen3.5-35B-A3B~\cite{PLACEHOLDER_qwen35_verify})
using free-form natural language.
An AI agent reads this review and selects rendering parameter adjustments
to address identified quality issues.
This loop repeats until the reviewer explicitly confirms the result is acceptable.
Section~\ref{sec:method:evolution} describes this stage in detail.

\begin{figure}[t]
  \centering
  % TODO: Insert pipeline overview figure
  \fbox{\parbox{0.95\linewidth}{\centering\vspace{2cm}
    \textbf{[Figure 1: Pipeline Overview]}\\
    Horizontal flowchart: Text Expansion $\to$ T2I Generation $\to$ SAM3 Segmentation
    $\to$ 3D Reconstruction $\to$ Scene-Aware Rendering $\to$ VLM Review Loop
    \vspace{2cm}}}
  \caption{
    Overview of the \aris{} pipeline.
    Starting from a natural-language object description,
    the pipeline automatically generates a high-quality RGBA image,
    reconstructs a textured 3D mesh, renders it in a real Blender scene,
    and iteratively refines rendering parameters through free-form VLM feedback.
  }
  \label{fig:pipeline}
\end{figure}


\subsection{Scene-Aware Blender Rendering}
\label{sec:method:blender}

Rendering each 3D object requires placing it plausibly within a scene,
matching the scene's lighting environment, and rendering from consistent viewpoints.
We use a fixed Blender scene file (a furnished indoor scene) as the background environment.

\paragraph{Scene environment preservation.}
A key design principle is \emph{not} to overwrite the scene's existing lighting:
the world environment (HDRI) and existing lights are preserved from the scene file.
Key lights are scaled by a fixed factor (0.3) rather than disabled,
ensuring the object receives natural scene illumination.
This avoids the common failure mode of rendering objects under artificial studio lighting
that clashes with the scene background.

\paragraph{Object placement.}
The object mesh is placed on the scene ground plane using a raycast procedure:
the system casts a ray downward from the object's bounding box to detect
the real ground surface height, ensuring physical plausibility.
The object is positioned such that its base rests on the scene floor,
with a configurable vertical offset for fine adjustment.

\paragraph{Rendering configuration.}
All renders use Cycles path tracing with 512 samples at $1024 \times 1024$ resolution,
providing high-quality global illumination and shadow detail.
For each object, we render eight horizontal viewpoints at
$0^\circ$, $45^\circ$, $90^\circ$, $135^\circ$, $180^\circ$, $225^\circ$, $270^\circ$, and $315^\circ$ azimuth,
corresponding to the natural-language view labels defined in Section~\ref{sec:dataset}.


\subsection{Free-Form VLM-Guided Rendering Evolution}
\label{sec:method:evolution}

The core contribution of \aris{} is the iterative quality refinement loop
that drives rendered images toward acceptable visual quality through VLM feedback.

\paragraph{VLM review.}
After each rendering round, a VLM (Qwen3.5-35B-A3B, with extended thinking enabled)
reviews the rendered images and produces a free-form natural language assessment.
The reviewer is prompted to identify specific visual problems,
including: flat or low-contrast lighting (\texttt{flat\_low\_contrast}),
color inconsistency between object and scene (\texttt{color\_shift}),
scene lighting mismatch (\texttt{scene\_light\_mismatch}),
missing or weak contact shadows (\texttt{shadow\_missing}),
physically implausible object placement (\texttt{floating}, \texttt{ground\_intersection}),
scale implausibility (\texttt{scale\_implausible}),
and overly smooth or plastic-looking materials (\texttt{material\_too\_plastic}).
The reviewer concludes with one of three verdicts:
\emph{keep} (result is acceptable), \emph{revise} (specific changes recommended),
or \emph{reject} (fundamental problems requiring re-rendering).

\paragraph{AI agent decision.}
An AI agent reads the full free-form reviewer text from the trace record
and identifies the primary diagnosed issue.
Based on the diagnosis, the agent selects one or more actions from the rendering action space
(Table~\ref{tab:action_space}).
This design differs from rigid score-based controllers:
the agent interprets semantic descriptions of visual failures
rather than responding mechanically to numeric thresholds.

\paragraph{Action space.}
The rendering action space contains 24 atomic actions organized into four groups:

\begin{table}[h]
\caption{Rendering action space. Each action adjusts one rendering parameter
by a fixed delta within enforced bounds.}
\label{tab:action_space}
\centering
\begin{tabular}{llcp{5.5cm}}
\toprule
Group & Example Actions & Count & Target Parameters \\
\midrule
Lighting & \texttt{L\_KEY\_UP}, \texttt{L\_KEY\_YAW\_POS\_15} & 4 &
  Key light scale, key light yaw angle \\
Object & \texttt{O\_LIFT\_SMALL}, \texttt{O\_SCALE\_UP\_10} & 6 &
  Vertical offset, yaw rotation, scale \\
Scene & \texttt{ENV\_ROTATE\_30}, \texttt{ENV\_STRENGTH\_UP} & 5 &
  HDRI yaw, environment strength, contact shadow \\
Material & \texttt{M\_VALUE\_DOWN}, \texttt{M\_ROUGHNESS\_UP} & 9 &
  Saturation, value, hue offset, roughness, specular \\
\midrule
\textbf{Total} & & \textbf{24} & \\
\bottomrule
\end{tabular}
\end{table}

\paragraph{Issue-to-action mapping.}
The agent uses a structured mapping from diagnosed issues to candidate actions.
For example, \texttt{flat\_low\_contrast} triggers environment strength adjustment (\texttt{ENV\_STRENGTH\_UP});
\texttt{color\_shift} triggers material value and saturation correction
(\texttt{M\_VALUE\_DOWN\_STRONG}, \texttt{M\_SATURATION\_DOWN});
\texttt{scene\_light\_mismatch} triggers environment and key light rotation
(\texttt{ENV\_ROTATE\_30}, \texttt{L\_KEY\_YAW\_POS\_15}).

\paragraph{Anti-oscillation and step decay.}
To prevent the agent from oscillating between opposing actions,
the system tracks the sign of each parameter's adjustment history.
If a parameter's sign has reversed three or more times, that parameter is frozen
for the remainder of the optimization.
Action step sizes decay with round number (100\% at round 0, 70\% at round 1, 50\% at round 2+)
to achieve fine-grained convergence near the optimum.

\paragraph{Termination.}
The loop terminates when the VLM reviewer explicitly produces a \emph{keep} verdict,
or after a maximum of five rounds.
Only results with an explicit keep verdict are considered quality-approved.

\begin{figure}[t]
  \centering
  % TODO: Insert VLM evolution loop detail figure
  \fbox{\parbox{0.95\linewidth}{\centering\vspace{2.5cm}
    \textbf{[Figure 2: VLM Evolution Loop]}\\
    Left column: Round 0 render $\to$ Round 1 render $\to$ $\cdots$ $\to$ Final (keep).\\
    Right column: VLM free-form text summary + agent-selected action at each round.
    \vspace{2.5cm}}}
  \caption{
    Illustration of the VLM-guided rendering evolution loop on a representative object.
    At each round, the VLM produces a free-form assessment of visual quality;
    the AI agent reads this text and selects a rendering action to address the primary issue.
    The loop continues until the reviewer issues a keep verdict.
  }
  \label{fig:evolution_loop}
\end{figure}
